\chapter{Euler Hamilton}
%%%%%%%%%%%%%%%%%%%%%%%%%%%%%%%%%%%%%%%%%%%%%%%%%%%%%%%%%%%%%%%%%%%%%%%%%%%%%%%
\section{Euler Hamilton}

% Generated by scripts/blueprint_restate.py from TCSlib.GraphTheory.Core.EulerHamilton.
% Each body is the STATEMENT only, taken from the declaration's Lean docstring:
% the one-line summary plus the verbatim book statement.  Proof, proof plan,
% significance commentary and formalisation notes are excluded by construction.

\begin{definition}[$G$ is eulerian: it has a closed Euler trail (an Euler tour)]
\lean{SimpleGraph.IsEulerianGraph}
\label{SimpleGraph.IsEulerianGraph}
\leanok
\texttt{G} is \textbf{eulerian}: it has a closed Euler trail (an Euler tour).

A trail that traverses every edge of $G$ is called an \emph{Euler trail} of $G$
because Euler was the first to investigate the existence of such trails in
graphs.

A \emph{tour} of $G$ is a closed walk that traverses each edge of $G$ at least once. An
\emph{Euler tour} is a tour which traverses each edge exactly once (in other words, a
closed Euler trail). A graph is \emph{eulerian} if it contains an Euler tour.
\end{definition}

\begin{definition}[Component count $\omega(G)$]
\lean{SimpleGraph.numComponents}
\label{SimpleGraph.numComponents}
\leanok
Component count \texttt{$\omega$(G)}.

\texttt{$\omega$(G)} is the number of components of \texttt{G} --- the classes into which the equivalence relation "is connected to" partitions the vertex set. Defined in chapter 1, not chapter 4, so there is no chapter-4 text to quote.
\end{definition}

\begin{definition}[join]
\lean{SimpleGraph.join}
\label{SimpleGraph.join}
\leanok
\textbf{Join} \texttt{G $\lor$ H}, rebuilt on Mathlib: within-side edges from
\texttt{$\oplus$g}, all cross edges from \texttt{completeBipartiteGraph}.
\texttt{symm}/\texttt{loopless} are inherited from \texttt{$\sqcup$}.

We first introduce the notion of the join of two graphs. The \emph{join} $G \vee H$ of
disjoint graphs $G$ and $H$ is the graph obtained from $G + H$ by joining each vertex of
$G$ to each vertex of $H$; it is represented diagrammatically as in figure 4.8.
\end{definition}

\begin{definition}[$C_{m,n} = K_m \lor (K_m^c + K_{n-2m})$, with $+$ the Mathlib disjoint sum $\oplus g$]
\lean{SimpleGraph.Cmn}
\label{SimpleGraph.Cmn}
\leanok
\uses{SimpleGraph.join}
\texttt{C\_\{m,n\} = K\_m $\lor$ (K\_m\textasciicircum{}c + K\_\{n-2m\})}, with \texttt{+} the Mathlib disjoint sum \texttt{$\oplus$g}.

Now, for $1 \le m < n/2$, let $C_{m,n}$ denote the graph
$K_m \vee (K_m^c + K_{n-2m})$, depicted in figure 4.9a; two specific examples,
$C_{1,5}$ and $C_{2,5}$, are shown in figures 4.9b and 4.9c.

That $C_{m,n}$ is nonhamiltonian follows immediately from theorem 4.2; for if
$S$ denotes the set of $m$ vertices of degree $n-1$ in $C_{m,n}$, we have
$\omega(C_{m,n} - S) = m + 1 > |S|$.

Theorem 4.6 then shows these are the \emph{extreme} nonhamiltonian graphs: every
nonhamiltonian simple graph is degree-majorised by some \texttt{C\_\{m,$\nu$\}}.
\texttt{C\_\{1,$\nu$\}} and \texttt{C\_\{2,5\}} reappear as the extremal cases of
Corollary 4.6.
\end{definition}

\begin{definition}[The sorted (ascending) degree sequence of $G$]
\lean{SimpleGraph.degreeSequence}
\label{SimpleGraph.degreeSequence}
\leanok
The sorted (ascending) \textbf{degree sequence} of \texttt{G}.

Listing the vertices as \texttt{v$_{1}$, \dots{}, v\_$\nu$}, the sequence
\texttt{(d(v$_{1}$), \dots{}, d(v\_$\nu$))} is a degree sequence of \texttt{G}. Chapter
4 always takes it in nondecreasing order --- as in the hypothesis of Theorem 4.5
(\S{}4.2, p. 65):

Let $G$ be a simple graph with degree sequence $(d_1, d_2, \ldots, d_\nu)$,
where $d_1 \leq d_2 \leq \ldots \leq d_\nu$ and $\nu \geq 3$.

which is why this definition sorts ascending.
\end{definition}

\begin{definition}[$G$ is degree-majorised by $H$: same order, and every sorted-degree entry of $G$ is $\le$ the\dots]
\lean{SimpleGraph.DegreeMajorised}
\label{SimpleGraph.DegreeMajorised}
\leanok
\uses{SimpleGraph.degreeSequence}
\texttt{G} is \textbf{degree-majorised} by \texttt{H}: same order, and every sorted-degree entry of \texttt{G}
is $\le$ the corresponding entry of \texttt{H}.

A sequence of real numbers $(p_1, p_2, \ldots, p_n)$ is said to be \emph{majorised} by
another such sequence $(q_1, q_2, \ldots, q_n)$ if $p_i \le q_i$ for $1 \le i \le n$. A
graph $G$ is \emph{degree-majorised} by a graph $H$ if $\nu(G) = \nu(H)$ and the
nondecreasing degree sequence of $G$ is majorised by that of $H$. For instance, the
5-cycle is degree-majorised by $K_{2,3}$ because $(2, 2, 2, 2, 2)$ is majorised by $(2,
2, 2, 3, 3)$.
\end{definition}

\begin{definition}[One Bondy--Chv\'atal closure step: add a nonadjacent pair whose degree sum is $\ge$ $\nu$]
\lean{SimpleGraph.ClosureStep}
\label{SimpleGraph.ClosureStep}
\leanok
One \textbf{Bondy--Chv\'atal closure step}: add a nonadjacent pair whose degree sum is
$\ge$ \texttt{$\nu$}.

Lemma 4.4.1 motivates the following definition. The \emph{closure} of $G$ is the graph
obtained from $G$ by recursively joining pairs of nonadjacent vertices whose degree sum
is at least $\nu$ until no such pair remains. We denote the closure of $G$ by $c(G)$.

$c(G)$ is well defined.

Its proof (the "first edge not in \texttt{G$_{2}$}" argument) is quoted on
\texttt{hamiltonian\_iff\_of\_closureSteps} below, which is where the content is used.
\end{definition}

\begin{definition}[$G$ is Hamilton-connected: every ordered pair is joined by a Hamilton path (Ex 4.2.11)]
\lean{SimpleGraph.IsHamiltonConnected}
\label{SimpleGraph.IsHamiltonConnected}
\leanok
\texttt{G} is \textbf{Hamilton-connected}: every ordered pair is joined by a Hamilton path (Ex 4.2.11).

\textbf{4.2.11} $G$ is \emph{Hamilton-connected} if every two vertices of $G$ are
connected by a Hamilton path. (a) Show that if $G$ is Hamilton-connected and $\nu \ge
4$, then $\varepsilon \ge [\frac{1}{2}(3\nu+1)]$. (b)* For $\nu \ge 4$, construct a
Hamilton-connected graph $G$ with $\varepsilon = [\frac{1}{2}(3\nu+1)]$. (J. W. Moon)
\end{definition}

\begin{definition}[$G$ is hypohamiltonian: not Hamiltonian, but every vertex-deleted subgraph is (Ex 4.2.12)]
\lean{SimpleGraph.IsHypohamiltonian}
\label{SimpleGraph.IsHypohamiltonian}
\leanok
\texttt{G} is \textbf{hypohamiltonian}: not Hamiltonian, but every vertex-deleted subgraph is (Ex 4.2.12).

\textbf{4.2.12} $G$ is \emph{hypohamiltonian} if $G$ is not hamiltonian but $G - v$ is
hamiltonian for every $v \in V$. Show that the Petersen graph (figure 4.4) is
hypohamiltonian. (Herz, Duby and Vigu\'e, 1967 have shown that it is, in fact, the
smallest such graph.)
\end{definition}

\begin{theorem}[Euler's theorem: eulerian graphs are those with all even degrees]
\lean{SimpleGraph.euler_tour_iff_no_odd_degree}
\label{SimpleGraph.euler_tour_iff_no_odd_degree}
\uses{SimpleGraph.IsEulerianGraph}
Let $G$ be a finite simple graph that is connected and has at least one edge. Then $G$
is eulerian — that is, it admits an Euler tour, a closed walk traversing each edge
exactly once — if and only if every vertex of $G$ has even degree.
\end{theorem}

\begin{theorem}[Euler trail criterion via odd-degree vertices]
\lean{SimpleGraph.euler_trail_iff_le_two_odd}
\label{SimpleGraph.euler_trail_iff_le_two_odd}
Let $G$ be a connected simple graph on a finite vertex set $V$. Then $G$ admits an
Eulerian trail — a walk between some pair of vertices $u,v$ that traverses every edge of
$G$ exactly once — if and only if the number of vertices of odd degree is at most $2$.
\end{theorem}

\begin{theorem}[Hamiltonian graphs are 1-tough]
\lean{SimpleGraph.hamiltonian_toughness}
\label{SimpleGraph.hamiltonian_toughness}
Let $G$ be a finite simple graph with vertex set $V$, and suppose $G$ has a Hamiltonian
cycle. Then for every nonempty proper subset $S \subsetneq V$, the number of connected
components of the induced subgraph $G - S$ obtained by deleting the vertices of $S$
satisfies
\[ c(G - S) \le \abs{S}. \]
\end{theorem}

\begin{theorem}[Dirac's theorem on Hamiltonian cycles]
\lean{SimpleGraph.dirac_hamiltonian}
\label{SimpleGraph.dirac_hamiltonian}
Let $G$ be a finite simple graph on $n$ vertices, where $n \ge 3$, and suppose every
vertex has degree at least $n/2$ (equivalently, $n \le 2\delta$, where $\delta$ denotes
the minimum degree of $G$). Then $G$ is Hamiltonian, that is, $G$ contains a cycle
passing through every vertex exactly once.
\end{theorem}

\begin{theorem}[Bondy–Chvátal edge-closure lemma for Hamiltonicity]
\lean{SimpleGraph.bondy_chvatal}
\label{SimpleGraph.bondy_chvatal}
Let $G$ be a finite simple graph on a vertex set with $n$ vertices, and let $u$ and $v$
be two distinct, nonadjacent vertices of $G$ whose degrees satisfy $\deg(u) + \deg(v)
\ge n$. Then $G$ is Hamiltonian if and only if the graph $G + uv$, obtained from $G$ by
adding the edge joining $u$ and $v$, is Hamiltonian.
\end{theorem}

\begin{theorem}[Closure steps preserve Hamiltonicity]
\lean{SimpleGraph.hamiltonian_iff_of_closureSteps}
\label{SimpleGraph.hamiltonian_iff_of_closureSteps}
\uses{SimpleGraph.ClosureStep}
Let $G$ be a finite simple graph on $\nu$ vertices, and suppose $H$ is obtained from $G$
by a finite sequence (possibly empty) of Bondy–Chvátal closure steps, each of which
selects two distinct nonadjacent vertices whose degrees sum to at least $\nu$ and joins
them by an edge. Then $G$ is Hamiltonian if and only if $H$ is Hamiltonian.
\end{theorem}

\begin{theorem}[Reaching the complete graph by closure steps forces a Hamiltonian cycle]
\lean{SimpleGraph.hamiltonian_of_closureSteps_top}
\label{SimpleGraph.hamiltonian_of_closureSteps_top}
\uses{SimpleGraph.ClosureStep}
Let $G$ be a finite simple graph on $\nu \ge 3$ vertices. Suppose the complete graph on
the vertices of $G$ can be obtained from $G$ by finitely many (possibly zero)
Bondy--Chv\'atal closure steps, each of which joins some pair of distinct nonadjacent
vertices whose degrees sum to at least $\nu$. Then $G$ is Hamiltonian, that is, $G$ has
a cycle through all of its vertices.
\end{theorem}

\begin{theorem}[Complete graphs on at least three vertices are Hamiltonian]
\lean{SimpleGraph.top_isHamiltonian}
\label{SimpleGraph.top_isHamiltonian}
Let $V$ be a finite vertex set with $\abs{V} \ge 3$. Then the complete graph on $V$ —
the simple graph in which every pair of distinct vertices is joined by an edge — is
Hamiltonian; that is, it contains a cycle passing through every vertex of $V$ exactly
once.
\end{theorem}

\begin{theorem}[Chvátal's degree condition for hamiltonicity]
\lean{SimpleGraph.chvatal_hamiltonian}
\label{SimpleGraph.chvatal_hamiltonian}
Let $G$ be a finite simple graph on $n \ge 3$ vertices, and write $\deg(v)$ for the
degree of a vertex $v$. Suppose that for every integer $m$ with $1 \le m$ and $2m < n$,
it is not the case that both $G$ has at least $m$ vertices of degree at most $m$ and $G$
has at least $n - m$ vertices of degree strictly less than $n - m$. Then $G$ is
Hamiltonian, that is, $G$ possesses a Hamiltonian cycle.
\end{theorem}

\begin{theorem}[Chvátal's degree-majorisation of nonhamiltonian graphs]
\lean{SimpleGraph.chvatal_degree_majorised}
\label{SimpleGraph.chvatal_degree_majorised}
\uses{SimpleGraph.Cmn, SimpleGraph.DegreeMajorised}
Let $G$ be a finite simple graph on $n \ge 3$ vertices, and suppose $G$ is not
Hamiltonian. Then there is an integer $m$ with $0 < m$ and $2m < n$ such that $G$ is
degree-majorised by $C_{m,n}$ (Chvátal, 1972).
\end{theorem}

\begin{theorem}[Edge-count bound forcing a Hamiltonian cycle]
\lean{SimpleGraph.ore_bondy_edge_bound}
\label{SimpleGraph.ore_bondy_edge_bound}
Let $G$ be a finite simple graph on $n$ vertices with $n \ge 3$. If the number of edges
of $G$ exceeds $\binom{n-1}{2} + 1$, then $G$ is Hamiltonian, that is, $G$ contains a
spanning cycle.
\end{theorem}

\begin{theorem}[Extremal non-Hamiltonian graphs at one more than the critical edge count]
\lean{SimpleGraph.ore_bondy_extremal}
\label{SimpleGraph.ore_bondy_extremal}
\uses{SimpleGraph.Cmn}
Let $G$ be a finite simple graph on $\nu \ge 3$ vertices that is not Hamiltonian and has
exactly $\binom{\nu-1}{2}+1$ edges. Then $G$ is isomorphic to $C_{1,\nu}$, unless $\nu =
5$, in which case $G$ is isomorphic to either $C_{1,5}$ or $C_{2,5}$.
\end{theorem}

\begin{theorem}[Cycle decomposition of an even graph]
\lean{SimpleGraph.even_degree_cycle_decomposition}
\label{SimpleGraph.even_degree_cycle_decomposition}
Let $G$ be a finite simple graph in which every vertex has even degree. Then the edges
of $G$ can be decomposed into cycles: there is a finite collection of cycles $C_1, C_2,
\ldots, C_m$ in $G$ such that every edge of $G$ lies in exactly one of the $C_i$.
\end{theorem}

\begin{theorem}[Edge-partition of a connected graph into trails]
\lean{SimpleGraph.odd_vertices_trail_cover}
\label{SimpleGraph.odd_vertices_trail_cover}
Let $G$ be a connected simple graph on a finite vertex set, and let $k \ge 1$. Suppose
that $G$ has exactly $2k$ vertices of odd degree. Then there exist $k$ trails $Q_1,
\dots, Q_k$ in $G$ such that every edge of $G$ lies in exactly one of them; that is, the
trails are pairwise edge-disjoint and their edge sets together exhaust $E(G)$.
\end{theorem}

\begin{theorem}[Graphs that are not 2-connected are nonhamiltonian]
\lean{SimpleGraph.nonhamiltonian_of_not_two_connected}
\label{SimpleGraph.nonhamiltonian_of_not_two_connected}
\uses{SimpleGraph.vertexConnectivity}
Let $G$ be a finite simple graph on at least three vertices, and suppose its vertex
connectivity satisfies $\kappa(G) < 2$ (equivalently, $G$ is not $2$-connected). Then
$G$ is nonhamiltonian; that is, $G$ has no Hamiltonian cycle.
\end{theorem}

\begin{theorem}[Non-Hamiltonicity of unbalanced bipartite graphs]
\lean{SimpleGraph.nonhamiltonian_of_unbalanced_bipartite}
\label{SimpleGraph.nonhamiltonian_of_unbalanced_bipartite}
Let $G$ be a finite simple graph, and let $X$ and $Y$ be disjoint sets of its vertices
such that every edge of $G$ joins a vertex of $X$ to a vertex of $Y$. If $\abs{X} \neq
\abs{Y}$, then $G$ is not Hamiltonian; that is, $G$ has no cycle passing through every
vertex of $G$.
\end{theorem}

\begin{theorem}[A Hamilton path bounds the components of a vertex-deletion]
\lean{SimpleGraph.hamiltonian_path_toughness}
\label{SimpleGraph.hamiltonian_path_toughness}
Let $G$ be a finite simple graph on vertex set $V$, and suppose $G$ has a Hamilton path,
that is, a path passing through every vertex of $G$ exactly once. Then for every
nonempty proper subset $S \subsetneq V$, the graph $G - S$ obtained by deleting the
vertices of $S$ has at most $\abs{S} + 1$ connected components.
\end{theorem}

\begin{theorem}[Chvátal's degree condition for a Hamilton path]
\lean{SimpleGraph.chvatal_hamiltonian_path}
\label{SimpleGraph.chvatal_hamiltonian_path}
Let $G$ be a simple graph on $\nu \ge 2$ vertices, with degrees arranged in
nondecreasing order as $d_1 \le d_2 \le \cdots \le d_\nu$. Suppose there is no integer
$m$ with $1 \le m < (\nu+1)/2$ for which both $d_m < m$ and $d_{\nu-m+1} < \nu - m$.
Then $G$ has a Hamilton path, that is, a path that passes through every vertex of $G$
exactly once. (Chvátal)
\end{theorem}

\begin{theorem}[Self-complementary graphs have a Hamilton path]
\lean{SimpleGraph.hamiltonian_path_of_self_complementary}
\label{SimpleGraph.hamiltonian_path_of_self_complementary}
Let $G$ be a finite simple graph, and suppose $G$ is self-complementary, meaning there
is a graph isomorphism between $G$ and its complement $G^c$. Then $G$ has a Hamilton
path: there exist vertices $a$ and $b$ and a walk from $a$ to $b$ in $G$ that visits
every vertex of $G$ exactly once.
\end{theorem}

\begin{theorem}[Degree-domination criterion for a Hamiltonian path (Clapham Ex. 4.2.5(a))]
\lean{SimpleGraph.hamiltonian_path_of_degree_dominates_compl}
\label{SimpleGraph.hamiltonian_path_of_degree_dominates_compl}
Let $G$ be a finite simple graph on $n$ vertices with complement $G^{c}$, and arrange
the vertex degrees of each graph in nondecreasing order, writing $d_{0} \le d_{1} \le
\cdots \le d_{n-1}$ for the sorted degree sequence of $G$ and $d'_{0} \le d'_{1} \le
\cdots \le d'_{n-1}$ for that of $G^{c}$. If $d'_{m} \le d_{m}$ for every index $m$ with
$2m \le n$, then $G$ has a Hamiltonian path: there are vertices $a$ and $b$ and a walk
from $a$ to $b$ that visits every vertex of $G$ exactly once.
\end{theorem}

\begin{theorem}[Erdős's edge-count sufficient condition for Hamiltonicity]
\lean{SimpleGraph.erdos_hamiltonian}
\label{SimpleGraph.erdos_hamiltonian}
Let $G$ be a finite simple graph with $n$ vertices and $m$ edges, and let $\delta$
denote its minimum degree. If $6\delta \le n$ and
\[
\binom{n-\delta}{2} + \delta^2 < m,
\]
then $G$ is Hamiltonian, that is, $G$ contains a cycle passing through every vertex
exactly once. (P. Erdős)
\end{theorem}

\begin{theorem}[Long path in a connected graph of large order (Dirac)]
\lean{SimpleGraph.long_path_of_order_gt_two_minDegree}
\label{SimpleGraph.long_path_of_order_gt_two_minDegree}
Let $G$ be a connected simple graph on a finite vertex set, with minimum degree
$\delta(G)$ and order $n$. If $2\,\delta(G) < n$, then $G$ contains a path of length at
least $2\,\delta(G)$; that is, there exist vertices $a$ and $b$ and a path from $a$ to
$b$ in $G$ with at least $2\,\delta(G)$ edges.
\end{theorem}

\begin{theorem}[A $2k$-regular graph on $4k+1$ vertices is Hamiltonian]
\lean{SimpleGraph.nash_williams_regular_hamiltonian}
\label{SimpleGraph.nash_williams_regular_hamiltonian}
Let $k \ge 1$ be an integer, and let $G$ be a finite simple graph on $4k+1$ vertices in
which every vertex has degree exactly $2k$. Then $G$ has a Hamiltonian cycle, that is, a
cycle passing through every vertex of $G$ exactly once.
\end{theorem}

\begin{theorem}[Edge bound for Hamilton-connected graphs]
\lean{SimpleGraph.hamiltonConnected_edge_bound}
\label{SimpleGraph.hamiltonConnected_edge_bound}
\uses{SimpleGraph.IsHamiltonConnected}
Let $G$ be a finite simple graph on $\nu \ge 4$ vertices with $\varepsilon$ edges. If
$G$ is Hamilton-connected — that is, every two distinct vertices of $G$ are joined by a
Hamilton path — then \[3\nu + 1 \le 2\varepsilon.\]
\end{theorem}
